\section{Model Validation and Testing}\label{sec:testing}

Beyond the per-question validations embedded in Section~\ref{sec:q1}--\ref{sec:q4},
we ran a cross-cutting test suite so that every number quoted in this paper is either
analytically derived, Monte-Carlo reproduced, or simulator-replayed from a fixed seed.

\subsection{Engine-Level Regression Tests}

The shared geometry engine (three-constraint predicate + obstacle-aware shortest path +
gimbal kinetics) is covered by a 51-check regression suite (Table~\ref{tab:testsuite}),
all passing. The suite pins the exact target projections (to $10^{-4}\,\mathrm{m}$),
exercises the closed-set boundary semantics (exactly-$1.5$-m clearance admissible,
tangent sightline blocked), and contains the dedicated sharp-corner regression
described in the table.

\begin{table}[!htb]\centering
\caption{Engine regression suite (51 checks, all passing).}
\label{tab:testsuite}
\begin{tabular}{lll}
\toprule
Module & What is pinned & Checks \\
\midrule
Target geometry & $T_1,T_2,T_3$ projections to $10^{-4}$ m & 3 \\
Clearance & field edge / facade touch / interior cases & 6 \\
Elevation & $r=r_{80}$ admissible, $r_{80}-0.1$ m rejected & 3 \\
Line of sight & through-roof, tangent face, $10^{-6}$ m off-face & 6 \\
Shortest path & direct $=$ Euclidean, detour bounds, per-segment gap & 9 \\
Sharp-corner regression & B1 SW corner straight segment admissible & 2 \\
Gimbal kinetics & shortest-arc slew, combined axes, reset & 6 \\
End-to-end & station feasibility + integrated schedule smoke & 16 \\
\bottomrule
\end{tabular}
\end{table}

Route legality is additionally certified by an independent segment-level auditor that
re-derives, for every leg of every final plan, the exact minimum distance to the field
boundary and to each building---e.g.\ the Q2(2) doubling vehicle's two legs at minimum
gaps of $10.0/10.0/14.0\,\mathrm{m}$ against the $1.5\,\mathrm{m}$ requirement.

\subsection{Statistical Validation Summary}

The statistical evidence is organized in four layers: \textbf{(L1)} analytic
identities (closed forms evaluated independently), \textbf{(L2)} Monte-Carlo
reproduction with sampling intervals, \textbf{(L3)} independent-implementation
cross-checks (Section~\ref{sec:testing} last subsection), and \textbf{(L4)} sensitivity
sweeps (Section~\ref{sec:sens}). Table~\ref{tab:mcsummary} maps each check to its
layer; every random experiment uses the fixed seed 2026.

\begin{table}[!htb]\centering
\caption{Monte-Carlo and refinement validations across the paper (fixed seed 2026;
layer tags in parentheses).}
\label{tab:mcsummary}
\begin{tabular}{llll}
\toprule
Where & Method & Metric & Result \\
\midrule
Q1 (\S\ref{sec:q1}) & $0.02$ m boundary refinement (L1) & area drift & $\le0.23\%$ \\
Q2(1) (\S\ref{sec:q21}) & exact-geometry MC, $2\times10^{5}$ (L2) & RMSE rel.\ diff & $0.0\%$ \\
Q2(1) (\S\ref{sec:q21}) & $\sigma_b\in\{0.5^\circ,1^\circ,2^\circ\}$ (L4) & optimal pair & invariant \\
Q2(2) (\S\ref{sec:q22}) & simulator replay (L3) & all starts & match \\
Q3 (\S\ref{sec:q3}) & exact-geometry MC, $10^{5}$/point (L2) & max $|\Delta P|$ & $0.18$ pp \\
Q3 (\S\ref{sec:q3}) & quadrature crossing vs.\ analytic (L1) & $s_{\max}$ & conservative $0.17$ m \\
Q4 (\S\ref{sec:q4}) & replay from captured states (L3) & all starts & match \\
\bottomrule
\end{tabular}
\end{table}

Table~\ref{tab:mcsummary} collects the statistical evidence: wherever an analytic curve
exists (Q2(1) RMSE, Q3 reliability), the exact-geometry Monte Carlo agrees to within
sampling error, and all random experiments use the fixed seed 2026. Sampling intervals
are stated analytically rather than by re-running: the $2\times10^{5}$-sample RMSE
estimate carries a standard error of $\mathrm{RMSE}/\sqrt{2N}=2.3\times10^{-4}\,\mathrm{m}$
($0.16\%$), so the observed $0.007\%$ relative deviation is well inside one standard
error; the Q3 per-point binomial standard errors are
$\sqrt{p(1-p)/10^{5}}\le7.0\times10^{-4}$ ($0.07$ pp), consistent with the Wilson
interval quoted at $s_{\max}$.

\subsection{Optimality Evidence and Honest Gaps}

For Q2(1), the Fisher-matrix lower bound is \emph{attained}, leaving no optimality gap.
For Q2(2), the per-alert lower bound gives $\max_a \mathrm{LB}_a=3.873\,\mathrm{s}$
against an achieved $11.046\,\mathrm{s}$: we report the honest $7.17\,\mathrm{s}$ gap
($6.78\,\mathrm{s}$ against the cross-implementation incumbent, a certified $3.6\%$
station-level residual, Section~\ref{sec:q22}) rather than claiming global optimality,
and the conservative incumbent is retained as the frozen Q4 baseline. For Q4, P1 is
exact within the same restricted search space, and P2's $\le4.53\,\mathrm{s}$ deficit
is certified by comparing both against identical replayed schedules.

\subsection{Cross-Implementation Reproducibility}

All headline numbers were re-derived by a fully independent second implementation
(different codebase, rounded-corner routing with an exact $R_V$-neighbourhood
clearance predicate, polar-candidate stations, global-objective coordinate descent;
cf.\ Section~\ref{sec:q22}, Validation~(4)). Results: Q1's optimal station
$(80.5004,63.6484)$, Q2(1)'s attaining pair V3+V5 (RMSE $0.14362\,\mathrm{m}$, travel
$52.47\,\mathrm{m}$) and Q3's $s_{\max}=66.92\,\mathrm{m}$ reproduce exactly; Q2(2)'s
incumbent improves to $10.651\,\mathrm{s}$ on the same structure family. The two
routers agree on every audited leg, and the refined schedule replays bit-compatibly in
this paper's own simulator---mutual oracles, not shared-derivative duplicates.

\section{Sensitivity Analysis}\label{sec:sens}

All perturbations below follow the one-factor-at-a-time plus region-map practice of
global sensitivity analysis \cite{saltelli2008}.

\subsection{Kinematic and Sensor Perturbations (Q2(2))}

The Q2(2) optimum was re-solved under $\pm10\%$ perturbations of the two dominant
kinematic parameters (Table~\ref{tab:sens}). The response is mild and monotone:
slower slewing ($-10\%$ gimbal rate) raises the worst alert start by
$+4.2\%$ ($11.046\to11.509\,\mathrm{s}$), faster slewing lowers it by $-3.4\%$;
$\pm10\%$ speed moves it by $+3.6\%/-1.5\%$. In all four cases the optimal
\emph{structure} is unchanged---only station positions and timings shift---so the
qualitative scheduling conclusion (double V3 on $T_1\!\to\!T_3$ with a shared station)
is robust to $\pm10\%$ parameter error.

\begin{table}[!htb]\centering
\caption{Q2(2) sensitivity: $\pm10\%$ in gimbal slew rate and vehicle speed.}
\label{tab:sens}
\begin{tabular}{lcccc}
\toprule
Case & max $t_{\mathrm{start}}$ (s) & $\Delta$ & makespan (s) & $\Delta$ \\
\midrule
Baseline                & 11.046 & ---    & 21.705 & ---    \\
Gimbal rate $-10\%$     & 11.509 & $+4.2\%$ & 22.852 & $+5.3\%$ \\
Gimbal rate $+10\%$     & 10.668 & $-3.4\%$ & 20.766 & $-4.3\%$ \\
Speed $-10\%$           & 11.439 & $+3.6\%$ & 22.414 & $+3.3\%$ \\
Speed $+10\%$           & 10.875 & $-1.5\%$ & 21.125 & $-2.7\%$ \\
\bottomrule
\end{tabular}
\end{table}

\subsection{Sensor-Noise and Geometry Robustness}

\textbf{Bearing-noise level (Q2(1)).} Re-running the pair optimization for
$\sigma_b\in\{0.5^\circ,1^\circ,2^\circ\}$ changes only the scale of the RMSE
($0.07181/0.14362/0.28725\,\mathrm{m}$, exactly proportional): the optimal pair set,
the $\gamma=90^\circ$ station geometry, and the travel tie-break are invariant. The
deployment recommendation does not depend on knowing $\sigma_b$ precisely.

\textbf{Boundary margins.} Q1's optimal stations sit at arc endpoints where either
clearance or line-of-sight binds; all constraints are evaluated as \emph{closed} sets
(touching allowed, entering forbidden), so the reported points are feasible exactly as
stated, and the $0.1$-m grid plus $0.02$-m refinement bounds the possible drift at
$0.23\%$ of area.

\textbf{Error-growth law (Q3).} The $(a,b)$ region of Fig.~\ref{fig:q3ab} is the
\emph{full} robustness map at the 80-m operating point: any $(a,b)$ with
$a+80b\le1.6384$ (intercept up to $1.64\,\mathrm{m}$ or slope up to
$0.0205\,\mathrm{m/m}$) sustains $95\%$ reliability, while the baseline $(0.3,0.02)$
needs $a+80b=1.9$, a $16\%$ budget overrun---a quantified, directional statement of how
far the stated sensor model is from compliance.

\textbf{Recovery-time sweep (Q4).} The $t_{rec}$ sweep in Fig.~\ref{fig:q4trec} acts as
a sensitivity analysis on the communication parameter: performance degrades at
$\approx0.49\,\mathrm{s}$ per second of outage (plateauing while vehicles idle), with no
threshold effect beyond the $t_{rec}\le10\,\mathrm{s}$ zero-gain band of pure family
(c).

\emph{Overall: every model's decision survives $\pm10\%$ kinematic perturbation and
wide noise-level variation; the only sharp sensitivities are directional ones (Q3's
error budget, Q4's outage duration), both of which are reported as explicit budgets
rather than hidden in point estimates.}
